\scanpage{14}\mspage{13}
one finds as \unsure{optimal|optimal} possibilities for defining the workshop structure on $\unsure{\mathcal{F}|\mathfrak{F}}$ :
\[
\left\{\begin{aligned}
 &\text{Either by}\quad
 \mathord{\tikz[baseline=(pair.base)]{\node[inner sep=5pt] (pair) {$\leq,\quad\ll$}; \draw ([xshift=1pt]pair.north west)--([xshift=-1pt]pair.north east); \draw ([yshift=-1pt]pair.north east)--([yshift=1pt]pair.south east); \draw ([xshift=-1pt]pair.south east)--(pair.south west)--([yshift=-2pt]pair.north west);}}\\
 &\text{or by}\quad\leq,\quad\preccurlyeq\\
 &\text{or by}\quad\ll,\quad
 \mathrel{\tikz[baseline=0pt,x=1em,y=1em]{\draw (-.38,.62)--(.28,.20)--(-.45,-.20); \draw (.35,.52)--(-.30,.18)--(.34,-.18); \draw (-.38,.62)--(-.40,.50)--(.28,.20); \draw (-.30,.18)--(-.42,.08)--(.32,-.14)--(.34,-.23);}}\ .
\end{aligned}\right.
\tag*{(1.28)}
\]
\margin{No, \unsure{its|that} will not suffice, \unsure{according to|according to} new point of view \ldots}

In the first version, I chose $\leq$, $\preccurlyeq$, defining in terms of these
\[
\left\{\begin{aligned}
 &\ll\ \text{is the preorder relation generated by}\\
 &\qquad\leq,\preccurlyeq,\quad \text{or again}\\
 &G\ll F\quad\Longleftrightarrow\quad\exists F'\ \text{with}\quad G\leq F'\preccurlyeq F
\end{aligned}\right.
\tag*{(1.29)}
\]
% slanted
\margin{I shall change the point of view on page 20, and shall no longer \unsure{assume} that $\ll$ is determined in this way in terms of \unsure{$\leq,\preccurlyeq$,|$\leq,\preccurlyeq$} but rather that $\preccurlyeq$ is determined by $\ll,\leq$}

This corresponds to the \deleted{\unsure{conceptual}} \inserted{\unsure{geometric|geometric}} and intuitive centre of gravity, which is indeed on the notion of subdivision, rather than on that of refinement. But technically it was a bit laborious, to deduce the properties of transitivity and \inserted{above all} of \unsure{reflexivity|reflexivity} for $\ll$, this led to putting in too many axioms. Still technically, it would seem that starting from $\leq,\ll$ is the simplest. This is the approach I shall follow shortly. As for the possibility of going via $\ll$, $\mathrel{\tikz[baseline=0pt,x=1em,y=1em]{\draw (-.38,.62)--(.38,.25)--(-.43,-.10); \draw (.34,.62)--(-.32,.25)--(.38,-.10);}}$, it is more of a curiosity than anything else,
